% TODO checklist:
% - [x] Inspected tests/ directory structure via list_dir
% - [x] Read tests/conftest.py for fixtures and configuration
% - [x] Inspected pyproject.toml for pytest markers and coverage config
% - [x] Read src/memgraph/test_mode.py for NL test mode
% - [x] Referenced Q&A.md Q44 for comprehensive testing details
% - [x] Inspected tests/fixtures/ for fake implementations
% - [x] Reviewed tests/e2e/playwright/ for E2E testing approach

\section{Testing Strategy}
\label{sec:testing}

The Cineca Agentic Platform employs a comprehensive multi-layer testing strategy to ensure correctness, reliability, and security across all platform components. This section documents the testing categories, infrastructure, and specialized testing approaches.

\subsection{Testing Categories and Markers}

The test suite is organized using pytest markers that classify tests by purpose, dependencies, and execution characteristics:

\begin{table}[ht]
\centering
\caption{Pytest markers and test categories.}
\label{tab:pytest-markers}
\small
\begin{tabular}{llll}
\hline
\textbf{Marker} & \textbf{Purpose} & \textbf{Speed} & \textbf{Dependencies} \\
\hline
\texttt{unit} & Pure Python logic & Fast ($<$1s each) & None \\
\texttt{integration} & Service interaction & Variable & Fakes or real services \\
\texttt{e2e} & Full workflow tests & Variable & App or live server \\
\texttt{security} & Auth/RBAC/rate limits & Variable & Auth fixtures \\
\texttt{performance} & Latency benchmarks & Slow & None (skipped by default) \\
\texttt{slow} & Long-running tests & Slow & Various \\
\texttt{memgraph\_nl} & NL$\rightarrow$Cypher smoke tests & $\sim$20 min & Memgraph, LLM \\
\texttt{memgraph\_nl\_full} & Full NL catalog tests & $\sim$90 min & Memgraph, LLM \\
\hline
\end{tabular}
\end{table}

\subsubsection{Running Tests by Category}

\begin{lstlisting}[language=bash, caption={Test execution commands by category.}]
# Full test suite
pytest -q

# Unit tests only (fast)
pytest -m "unit"

# Integration tests
pytest -m "integration"

# End-to-end tests
pytest -m "e2e"

# Security-focused tests
pytest -m "security"

# Performance tests (requires --runslow flag)
pytest -m "performance" --runslow

# Exclude slow tests for rapid iteration
pytest -m "not slow"
\end{lstlisting}

\subsection{Test Suite Statistics}

The test suite demonstrates comprehensive coverage of platform functionality:

\begin{table}[ht]
\centering
\caption{Test suite statistics.}
\label{tab:test-stats}
\small
\begin{tabular}{lr}
\hline
\textbf{Metric} & \textbf{Value} \\
\hline
Total test files & 272 \\
Test categories (directories) & 27 \\
Total lines of test code & $\sim$64,500 \\
Test functions & 2,700+ \\
Coverage threshold (minimum) & 60\% \\
Core module coverage target & 80\% \\
Router module coverage target & 70\% \\
\hline
\end{tabular}
\end{table}

\subsection{Testing Infrastructure Details}
\label{subsec:testing-infrastructure}

The testing infrastructure provides fixtures for isolated test execution and dependency mocking.

\subsubsection{Session-Scoped Event Loop}

Async tests share a session-scoped event loop for efficiency:

\begin{lstlisting}[language=Python, caption={Event loop fixture from \texttt{tests/conftest.py}.}]
import pytest
import asyncio

@pytest.fixture(scope="session")
def event_loop() -> Generator[asyncio.AbstractEventLoop, None, None]:
    """Session-scoped event loop for async tests."""
    loop = asyncio.new_event_loop()
    yield loop
    loop.close()
\end{lstlisting}

\subsubsection{FastAPI Test Client}

Both synchronous and asynchronous test clients are provided:

\begin{lstlisting}[language=Python, caption={Test client fixtures.}]
@pytest.fixture
def client(app):
    """Synchronous TestClient for simple tests."""
    from fastapi.testclient import TestClient
    with TestClient(app) as tc:
        yield tc

@pytest_asyncio.fixture
async def async_client(app) -> AsyncGenerator[Any, None]:
    """Async HTTP client with ASGI transport and lifespan."""
    import httpx
    
    transport = httpx.ASGITransport(app=app, lifespan="on")
    async with httpx.AsyncClient(
        transport=transport, 
        base_url="http://test"
    ) as ac:
        yield ac
\end{lstlisting}

\subsubsection{Settings Patching}

Tests can override configuration settings:

\begin{lstlisting}[language=Python, caption={Settings patch fixture.}]
@pytest.fixture
def settings_patch(monkeypatch) -> Callable[..., None]:
    """Patch src.config.settings for test duration."""
    from src.config import settings
    
    def _apply(**kwargs: Any) -> None:
        for k, v in kwargs.items():
            monkeypatch.setattr(settings, k, v, raising=False)
    
    return _apply

# Usage in tests:
def test_rate_limiting(settings_patch):
    settings_patch(
        RATE_LIMIT_ENABLED=True,
        RATE_LIMIT_DEFAULT_LIMIT=5
    )
    # Test rate limiting behavior...
\end{lstlisting}

\subsubsection{OIDC Token Fixtures}

Authentication tests use generated RSA keypairs and JWT tokens:

\begin{lstlisting}[language=Python, caption={OIDC fixtures for authentication testing.}]
@pytest.fixture(scope="session")
def oidc_keys(tmp_path_factory) -> Dict[str, Any]:
    """Generate test RSA keypair and JWKS file."""
    from tests.fixtures.oidc import generate_rsa_keypair, write_jwks
    
    pair = generate_rsa_keypair()
    jwks_dir = tmp_path_factory.mktemp("jwks")
    jwks_path = jwks_dir / "jwks.json"
    write_jwks(jwks_path, pair["public_jwk"])
    
    return {
        "kid": pair["kid"],
        "private_pem": pair["private_pem"],
        "public": pair["public_jwk"],
        "jwks_path": jwks_path,
        "issuer": "https://test-issuer.local/",
        "audience": "cineca-api",
    }
\end{lstlisting}

\subsection{Database Handling in Tests}

Tests use appropriate strategies for each database type to ensure isolation and reproducibility.

\subsubsection{PostgreSQL Testing}

PostgreSQL tests use session rollback for isolation:

\begin{lstlisting}[language=Python, caption={PostgreSQL test session fixture.}]
@pytest.fixture
def db_session():
    """Create isolated database session with rollback."""
    from db.postgres_control.database import SessionLocal
    
    session = SessionLocal()
    try:
        yield session
    finally:
        session.rollback()
        session.close()
\end{lstlisting}

\subsubsection{Redis Testing}

Redis tests can use either real Redis or an in-memory fake:

\begin{lstlisting}[language=Python, caption={Redis test fixture with fallback.}]
@pytest.fixture
def redis_fake():
    """In-memory Redis stub when REDIS_URL not reachable."""
    if _can_reach_redis():
        # Use real Redis with unique namespace prefix
        return RealRedisWithNamespace(prefix=f"test:{uuid.uuid4().hex}")
    return FakeRedis()

class FakeRedis:
    """Deterministic in-memory Redis double."""
    
    def __init__(self):
        self._data: dict[str, Any] = {}
        self._expiry: dict[str, float] = {}
    
    async def get(self, key: str) -> Any:
        if key in self._expiry and time.time() > self._expiry[key]:
            del self._data[key]
            del self._expiry[key]
            return None
        return self._data.get(key)
    
    async def set(self, key: str, value: Any, ex: int | None = None):
        self._data[key] = value
        if ex:
            self._expiry[key] = time.time() + ex
\end{lstlisting}

\subsubsection{Memgraph Testing}

Memgraph tests use a fake adapter for deterministic behavior:

\begin{lstlisting}[language=Python, caption={Fake Memgraph adapter from \texttt{tests/fixtures/fake\_memgraph.py}.}]
class FakeMemgraphAdapter:
    """Deterministic in-memory Memgraph double."""
    
    def __init__(self):
        self._nodes: dict[str, dict] = {}
        self._edges: list[dict] = []
    
    def execute(self, query: str, params: dict = None) -> list[dict]:
        """Parse Cypher-like queries and return mock results."""
        if "RETURN 1" in query:
            return [{"ok": 1}]
        
        # Pattern matching for common test queries
        if "MATCH (u:User)" in query:
            return self._mock_user_results(params)
        
        return []
    
    def _mock_user_results(self, params: dict) -> list[dict]:
        # Return deterministic test data
        return [{"name": "Test User", "id": "test-001"}]
\end{lstlisting}

\subsubsection{Switching to Real Services}

Tests can use real external services when available:

\begin{lstlisting}[language=bash, caption={Running tests with real services.}]
# Use real Memgraph
export MG_HOST=localhost MG_PORT=7687
pytest -m "integration" --run-real-memgraph

# Use real Redis
export REDIS_URL=redis://localhost:6379/0
pytest -m "integration"

# Use real Auth0 tokens
export AUTH0_DOMAIN=your-domain.auth0.com
pytest -m "e2e" --use-auth0
\end{lstlisting}

\subsection{Memgraph NL Test Mode}
\label{subsec:memgraph-nl-test-mode}

The platform provides a deterministic testing mode for the NL$\rightarrow$Cypher pipeline that eliminates LLM output variability.

\subsubsection{Test Mode Configuration}

Test mode is controlled via environment variables:

\begin{lstlisting}[language=Python, caption={NL test mode configuration from \texttt{src/memgraph/test\_mode.py}.}]
# Environment variables for NL test mode
_TEST_MODE_ENV = "LLM_MEMGRAPH_NL_TEST_MODE"
_PROMPT_PATH_ENV = "LLM_MEMGRAPH_NL_PROMPTS_PATH"
_DEFAULT_PROMPTS_REL = Path("tests/integration/resources/memgraph_nl_prompts.json")

def _is_enabled() -> bool:
    """Check if NL test mode is enabled."""
    value = os.getenv(_TEST_MODE_ENV)
    if value is None:
        # Default enabled for integration tests
        return True
    return value.strip().lower() in {"1", "true", "yes", "on"}
\end{lstlisting}

\subsubsection{Prompt Catalog}

Test prompts are defined in a JSON catalog that maps natural language to expected Cypher:

\begin{lstlisting}[language=json, caption={Sample NL prompt catalog entry.}]
[
  {
    "id": "user-count-001",
    "text": "How many users are in the system?",
    "category": "aggregation",
    "expected_cypher": "MATCH (u:User) RETURN count(u) AS total",
    "expected_result_shape": {"total": "number"},
    "todo_mode": "single_query"
  },
  {
    "id": "user-institution-002",
    "text": "Which users work at CINECA?",
    "category": "relationship",
    "expected_cypher": "MATCH (u:User)-[:WORKS_AT]->(i:Institution {name: 'CINECA'}) RETURN u.name",
    "expected_result_shape": {"u.name": "string[]"}
  }
]
\end{lstlisting}

\subsubsection{Hint Retrieval}

During test execution, prompts are matched to their expected outputs:

\begin{lstlisting}[language=Python, caption={Prompt hint retrieval function.}]
def get_prompt_hints(prompt: str | None) -> Optional[Dict[str, Any]]:
    """Return prompt metadata when NL test mode is enabled."""
    if not _is_enabled() or not prompt:
        return None
    
    prompt_path = os.getenv(_PROMPT_PATH_ENV)
    if not prompt_path:
        prompt_path = str(_default_prompt_path())
    
    index = _load_prompt_index(prompt_path)
    normalized = _normalize(prompt)
    
    entry = index.get(normalized)
    if entry:
        log.debug(
            "memgraph.prompts.match",
            prompt_preview=prompt[:80],
            prompt_id=entry.get("id"),
            category=entry.get("category"),
        )
    return entry
\end{lstlisting}

\subsubsection{Benefits of NL Test Mode}

\begin{enumerate}
    \item \textbf{Deterministic Results}: Tests produce consistent outputs independent of LLM variability
    \item \textbf{Fast Execution}: No LLM API calls required for matched prompts
    \item \textbf{Offline Testing}: Tests can run without network connectivity
    \item \textbf{Regression Detection}: Changes to prompt handling are immediately visible
    \item \textbf{Coverage Tracking}: Catalog documents supported query patterns
\end{enumerate}

\subsection{Continuous Integration and Quality Gates}
\label{subsec:ci-quality-gates}

The project enforces quality gates through automated checks in the CI pipeline.

\subsubsection{Test Execution in CI}

\begin{lstlisting}[language=bash, caption={CI test execution configuration.}]
# CI runs subset for speed, local runs full suite
pytest --maxfail=1 --disable-warnings -q \
    --cov=src --cov-report=term-missing:skip-covered

# Parallel execution for faster CI
pytest -n auto  # Requires pytest-xdist
\end{lstlisting}

\subsubsection{Coverage Requirements}

Coverage thresholds are enforced per module category:

\begin{lstlisting}[language=Python, caption={Coverage configuration from \texttt{pyproject.toml}.}]
[tool.coverage.run]
branch = true
source = ["src"]
omit = ["*/tests/*", "*/test_*.py", "*/__pycache__/*"]

[tool.coverage.report]
fail_under = 60  # Overall threshold

# Per-module thresholds (enforced via custom rules)
# Core modules: 80% (src/mcp/core/*, src/security/*)
# Router modules: 70% (src/routers/*)
# Service modules: 65% (src/services/*)

exclude_lines = [
    "pragma: no cover",
    "if TYPE_CHECKING:",
    "raise NotImplementedError",
    "@(abc\\.)?abstractmethod",
]
\end{lstlisting}

\subsubsection{Quality Gate Checks}

\begin{table}[ht]
\centering
\caption{CI quality gate checks.}
\label{tab:ci-gates}
\small
\begin{tabular}{lll}
\hline
\textbf{Check} & \textbf{Tool} & \textbf{Failure Condition} \\
\hline
Code formatting & Black & Any formatting differences \\
Linting & Ruff & Any linting errors \\
Type checking & mypy & Any type errors \\
Security scan & Bandit & High severity findings \\
Unit tests & pytest & Any test failures \\
Coverage & coverage.py & Below 60\% overall \\
\hline
\end{tabular}
\end{table}

\subsection{Test Examples by Category}

\subsubsection{Unit Test Example}

\begin{lstlisting}[language=Python, caption={Unit test for PII scrubbing.}]
@pytest.mark.unit
def test_scrub_email():
    """Test that email addresses are properly redacted."""
    from src.security.pii_scrubber import scrub
    
    text = "Contact jane.doe@example.com for help"
    cleaned, findings = scrub(text)
    
    assert "example.com" not in cleaned
    assert any(f["type"] == "EMAIL" for f in findings)
    assert "[EMAIL_REDACTED]" in cleaned
\end{lstlisting}

\subsubsection{Integration Test Example}

\begin{lstlisting}[language=Python, caption={Integration test for archive service.}]
@pytest.mark.integration
@pytest.mark.asyncio
async def test_archive_round_trip(tmp_backups_dir, memgraph_fake):
    """Test backup and restore cycle."""
    from src.services.archive import ArchiveService
    
    svc = ArchiveService(
        etl=memgraph_fake.etl, 
        base_dir=tmp_backups_dir
    )
    
    # Create snapshot
    snap = await svc.snapshot_graph()
    assert snap.ok
    assert snap.data["file"].endswith(".json.gz")
    
    # Restore from snapshot
    res = await svc.restore_graph(snap.data["file"])
    assert res.ok
\end{lstlisting}

\subsubsection{Security Test Example}

\begin{lstlisting}[language=Python, caption={Security test for role-based access.}]
@pytest.mark.security
def test_admin_endpoint_requires_admin_role(client, user_token):
    """Test that admin endpoints reject non-admin users."""
    response = client.get(
        "/v1/admin/stats",
        headers={"Authorization": f"Bearer {user_token}"}
    )
    
    assert response.status_code == 403
    assert "admin" in response.json().get("detail", "").lower()
\end{lstlisting}

\subsubsection{E2E Test Example}

\begin{lstlisting}[language=Python, caption={End-to-end health check test.}]
@pytest.mark.e2e
def test_health_endpoints_up(client):
    """Verify all health endpoints return 200."""
    endpoints = ["/health", "/health/live", "/health/ready"]
    
    for endpoint in endpoints:
        response = client.get(endpoint)
        assert response.status_code == 200, f"{endpoint} failed"
\end{lstlisting}

\subsection{Common Test Fixtures Summary}

\begin{table}[ht]
\centering
\caption{Common test fixtures and their purposes.}
\label{tab:common-fixtures}
\small
\begin{tabular}{ll}
\hline
\textbf{Fixture} & \textbf{Purpose} \\
\hline
\texttt{event\_loop} & Session-scoped async event loop \\
\texttt{app} & FastAPI application instance \\
\texttt{client} & Synchronous test client \\
\texttt{async\_client} & Async HTTPX client with lifespan \\
\texttt{db\_session} & PostgreSQL session with rollback \\
\texttt{redis\_fake} & In-memory Redis or namespaced real Redis \\
\texttt{memgraph\_fake} & Deterministic Memgraph adapter \\
\texttt{oidc\_keys} & RSA keypair and JWKS for auth testing \\
\texttt{settings\_patch} & Configuration override helper \\
\texttt{auth\_headers} & JWT tokens for different roles \\
\texttt{test\_tenant} & Isolated tenant for testing \\
\texttt{mock\_llm\_provider} & Stubbed LLM responses \\
\hline
\end{tabular}
\end{table}

